\section{Discussion}\label{sec:discussion}

\paragraph{What the comparison shows.} The embedding dominates every identity metric, but the interpretation is narrower than ``deep beats classical''. Map A's reported metrics use original full-dimensional pixels; its fitted PCA basis is for the plots and full-rank isometry check, not for verification. Map B is a 23.5-million-parameter network trained on 3.31 million images and then frozen. The gap therefore combines representation, training corpus, and architecture, which this study cannot separate. Fitting PCA on the probe images makes the visualisations transductive but cannot establish that the full-vector verification gap is conservative.

\paragraph{The metrics disagree, and that is the point.} Raw pixels look respectable on identification (94.5\% 1-NN) but poor on silhouette and substantially worse on known-population verification. Identification asks whether the single nearest gallery image is correct, and with ten images per identity and pre-aligned data the nearest genuine image is often close enough. Verification fixes one threshold and requires the genuine and impostor distributions to be separated, not merely ordered. Map A's overlap makes that impossible; Map B's near-disjoint distributions make it easy. A good nearest-neighbour result is not evidence that a usable deployment threshold exists.

\paragraph{The prettier map is the worse space.} Map A's 2-D plot is a more faithful projection of its own space ($\rho=0.807$) than Map B's plot is of the embedding space ($\rho=0.517$), and Map A's 2-D axes are literally brightness and a lateral lighting gradient. A viewer comparing the two scatters is not comparing recognition ability: the plots preserve unequal fractions of the geometry they summarise. Raw-pixel variance is dominated by a few illumination directions, so two components capture much of its structure even though that structure is not identity; the embedding's leading variance directions are less identity-discriminative, so compressing it discards more of what makes it useful.

\paragraph{Observed lighting-condition contrasts.} The two front-lit room-versus-flash groups differ by mean distance 0.258, within the 0.227--0.371 same-condition repeat range. Across four flash-direction labels the mean pair distance is 0.714, or 3.83 times within-scan jitter, with side-to-bottom pairs largest. These are separate sequential captures of one subject: pose and auto-exposure are unmeasured, so the study does not identify a causal intensity or direction effect. Left--right (0.511) is smaller than front--side pairs, consistent with a tentative ``half-shadowed'' region. Adding a side-lit enrolment scan is a hypothesis for a future held-out verification experiment, not a validated policy.

\paragraph{The security--convenience dial.} At matched 1\% FRR, Map A accepts 48.6\% of impostors and Map B 0.044\%; at 1\% FAR, Map B accepts 99.9\% of genuine probes against 74.4\% for Map A. Within the sampled operating points Map B has a broad region with both error rates small, while Map A has none. Thresholds do not transfer: the verification EER thresholds are 8.51 for pixels and 0.84 for embeddings (the pair-level EER thresholds are different again), and webcam distances live on a third scale. The cross-domain check makes the point: the smallest Olivetti-to-webcam distance (0.887) is below the worst same-person webcam distance (0.967), but the two numbers are not comparable scores, because the cross-domain value is confounded by device, crop, resolution, and population and is not a calibrated impostor score. It is not evidence that a stranger would be accepted, and it is neither an upper nor a lower bound on webcam-domain stranger distance. Domain shift is large enough to make a threshold fitted on one dataset meaningless on another; any deployment must calibrate on in-domain impostors, which this study cannot provide because it has no stranger capture. The same-sample operating points in Section~\ref{sec:results} are therefore exploratory, and an additive identity-disjoint check does not close the gap: at a nominal 1\% calibration FAR, Map B reaches held-out TAR 100\% against 75.08\% for Map A, but it still accepts 12.1\% of never-enrolled identities as somebody (any-of-13-gallery claims; Map A 18.9\%), and the nominal 1\% FAR is not validated out of sample (per-claim held-out FAR is 0.77\%). Unknown-identity rejection remains the open problem.
